\begin{exsol}{20}{\MR{} Food for Thought 5.1}{163--164}{Exam}{B2}
\asked Suppose people are partly ``hard-wired'' with a taste for vengeance and a conscience, a
preference for honesty and fairness (Frank's evolutionary argument). How would such preferences
change the chapter's analyses of commitment and bargaining?

\method Commitment and reneging, and the grab game of Table 5.1 (\MR{} pp.~133, 138--139);
reputation (pp.~139--140); bargaining with private information and the honest-reporting
constraints of Tables 5.2--5.3 (pp.~140--143); the ultimatum game and its experimental evidence
(Varian 2014, pp.~582--583). Formally, the preferences add a psychic term to material payoffs,
and we ask which constraints of the chapter it relaxes or tightens.

\solution
\textbf{1. Emotions are commitment devices.} The chapter's recurring problem is that a selfish
player who keeps an option after the fact cannot promise or threaten credibly. A conscience makes
\emph{promises} credible. In the grab game of Table 5.1 (p.~138) a selfish firm grabs whatever
the other does (3 > 2 and $-1 > -2$). If grabbing costs a guilt payoff $g$, refraining becomes
dominant when $g > 1$, the mutual-refrain outcome is an equilibrium, and the specific investment
that opportunism ``completely destroys'' is undertaken. A taste for vengeance makes
\emph{threats} credible. Punishing a cheater is costly ex post, which is why parents and
governments facing hostage-takers are tempted to forgive (box, p.~135); a player who
\emph{enjoys} retaliating no longer faces that temptation, so the threat deters, and in
equilibrium it seldom has to be carried out. Neither device burns a fleet or builds an
overspecialized plant (p.~139): the commitment is carried in the person.

\textbf{2. Bargaining predictions change.} With selfish players a one-shot offer is pure
commitment: the responder accepts anything positive and the proposer keeps almost all (``something
like 99 cents'', Varian p.~582). A responder known to reject offers below a threshold $k$ receives
$k$; vengeance is exactly the commitment to walk away that the selfish responder lacks. This is
what experiments show: average offers of about 45 cents, usually accepted, with small offers
rejected at a cost to the responder (Varian pp.~582--583). Proposers then behave nearly optimally
given the rejection rates, so selfishness plus correct expectations explains proposers, and
fairness explains responders. In alternating offers the same logic raises each player's
effective threat point, so the predicted split moves toward the norm (often 50--50), not toward
the more patient player. A further consequence: when people care about \emph{how} surplus is
divided, the value maximization principle loses its clean separation between the value-creating
decision and the transfers (\MR{} Ch.~2), because the transfer itself enters utility.

\textbf{3. Honesty relaxes the private-information constraints.} The impossibility on p.~142
comes from incentive-compatibility constraints that assume lying is free. If misreporting costs
$\psi$, the seller's constraint becomes $p \ge 1.75 - 1.25\psi$ and the buyer's
$p \le 1.25 + 1.25\psi$: a conscience worth $\psi \ge \$0.20$ restores efficient trade at
$p = \$1.50$. Rents paid for honesty can be partly replaced by honesty itself.

\textbf{4. But the same preferences also create costs.} (i) \emph{Disagreement.} When the
threshold $k$ is private, a proposer must trade a bigger offer against a higher acceptance
probability, so some offers are rejected and surplus is destroyed. This is the inefficiency of
Tables 5.2--5.3 with fairness types in place of values. (ii) \emph{Self-serving fairness.} \MR{}
warn that ``what is honest and fair is rarely as obvious in the rich settings of reality'' and
that people tend to see as fair what serves their own interests (pp.~138--139). Two vengeful
parties with different fairness standards become a strike or a feud, not a deal. (iii)
\emph{Credibility of the preference itself.} Frank's argument works only if the traits are
recognizable; otherwise the selfish can mimic them, and the commitment value collapses back into
a reputation problem (pp.~139--140). Hence \MR's caution: it is ``too risky to rely on others to
act consistently contrary to their own selfish interests unless there is something that commits
them''.

\textbf{Summary.} Such preferences enlarge what can be committed to (more investment, less
hold-up, honest reporting, offers near equal division) but add a new source of failure,
disagreement driven by conflicting or privately known notions of fairness.

\interp For the course: commitment does not need to be contractual, and emotions, norms and
corporate culture can carry it. An organization designer should treat them as a complement to
explicit incentives, not a substitute, because their strength is private information and their
content is self-serving. The ultimatum evidence is also a reminder that equilibrium predictions
depend on the payoffs actually in players' heads.

\verify The arithmetic is checked in \texttt{check\_D.py}: in Table 5.1, refraining is dominant
iff $g > 1$; with a lying cost $\psi$, \MR's bounds become $1.75 - 1.25\psi$ and
$1.25 + 1.25\psi$ and meet at $\psi = 0.2$, $p = 1.50$, the same price \MR{} use for the
reversed-probability case (p.~143). The experimental magnitudes are Varian's (pp.~582--583).

\answer{Fairness and vengeance act as built-in commitment devices. A conscience makes promises
credible and protects specific investments (guilt above 1 removes grabbing in Table 5.1); a
taste for revenge makes threats credible and gives responders bargaining power, which explains
offers of about 45 cents instead of the SPNE near zero. Honesty relaxes incentive-compatibility
constraints (a lying cost of \$0.20 restores efficient trade in \MR's example). The cost is that
privately known or self-serving notions of fairness bring new disagreement, and the preferences
help only if they are observable and cannot be mimicked.}
\end{exsol}
